\gdef\tagBrand{NT}
\gdef\footerAddress{Zalo: 0799666070}
\setcounter{ex}{0}

% --- CÂU 1 ---
\begin{ex}
	Khẳng định nào sau đây là một mệnh đề đúng?
	\choice
	{$15$ là số nguyên tố}
	{$\sqrt{7}$ là số hữu tỉ}
	{\True Số $0$ là số tự nhiên}
	{Phương trình $x^2 + 1 = 0$ có nghiệm thực}
	\loigiai{
		\begin{itemize}
			\item Phương án A sai vì $15 = 3 \times 5$ nên $15$ là hợp số.
			\item Phương án B sai vì $\sqrt{7}$ là số vô tỉ.
			\item Phương án C đúng vì $0 \in \mathbb{N}$.
			\item Phương án D sai vì $x^2 + 1 \ge 1 > 0$ với mọi $x \in \mathbb{R}$, nên phương trình vô nghiệm thực.
		\end{itemize}
	}
\end{ex}

% --- CÂU 2 ---
\begin{ex}
	Mệnh đề phủ định của mệnh đề $P\colon \text{``}\exists x \in \mathbb{R},\ x^2 - 4x + 4 \le 0\text{''}$ là
	\choice
	{\True $\overline{P}\colon \text{``}\forall x \in \mathbb{R},\ x^2 - 4x + 4 > 0\text{''}$}
	{$\overline{P}\colon \text{``}\exists x \in \mathbb{R},\ x^2 - 4x + 4 > 0\text{''}$}
	{$\overline{P}\colon \text{``}\forall x \in \mathbb{R},\ x^2 - 4x + 4 \ge 0\text{''}$}
	{$\overline{P}\colon \text{``}\forall x \notin \mathbb{R},\ x^2 - 4x + 4 > 0\text{''}$}
	\loigiai{
		Phủ định của kí hiệu $\exists$ là $\forall$, phủ định của bất đẳng thức $\le$ là $>$.\\
		Do đó, mệnh đề phủ định của $P$ là $\overline{P}\colon \text{``}\forall x \in \mathbb{R},\ x^2 - 4x + 4 > 0\text{''}$.
	}
\end{ex}

% --- CÂU 3 ---
\begin{ex}
	Tập hợp $E = \{x \in \mathbb{R} \mid -2 \le x < 3\}$ được biểu diễn dưới dạng khoảng, đoạn, nửa khoảng là
	\choice
	{$(-2; 3]$}
	{\True $[-2; 3)$}
	{$[-2; 3]$}
	{$(-2; 3)$}
	\loigiai{
		Do $x$ lớn hơn hoặc bằng $-2$ và nhỏ hơn $3$ nên ta sử dụng ngoặc vuông tại đầu mút $-2$ và ngoặc tròn tại đầu mút $3$.\\
		Vậy $E = [-2; 3)$.
	}
\end{ex}

% --- CÂU 4 ---
\begin{ex}
	Cho hai tập hợp $A = \{1; 2; 3; 4\}$ và $B = \{2; 4; 6\}$. Tập hợp $A \setminus B$ bằng
	\choice
	{\True $\{1; 3\}$}
	{$\{2; 4\}$}
	{$\{6\}$}
	{$\{1; 2; 3; 4; 6\}$}
	\loigiai{
		Tập hợp $A \setminus B$ gồm các phần tử thuộc $A$ nhưng không thuộc $B$.\\
		Vì $A = \{1; 2; 3; 4\}$ và các phần tử $2, 4 \in B$ nên $A \setminus B = \{1; 3\}$.
	}
\end{ex}
